\chapter{Constructive Status and Routes to Continuum QFT}
\label{ch:constructive-status-routes}

\ConventionsEuclidean

Constructive QFT asks for existence theorems for the objects used elsewhere in
the monograph: Schwinger functions, Wightman functions, Hilbert spaces, local
fields or local observable algebras, symmetry actions, positivity, and
continuum limits.  A regulator is part of a construction only after a limit is
specified and the limiting object is shown to satisfy the required axioms or
replacement framework.

\section{Three Constructive Routes}

\begin{definition}[Constructive realization]
\label{def:constructive-realization}
A constructive realization of a QFT in a specified framework consists of:
correlation functions or an observable net; positivity or reflection
positivity; covariance or the declared background-functorial replacement;
locality; a reconstruction theorem or direct Hilbert-space construction; and
control of the limiting procedure in a topology strong enough to verify those
properties.
\end{definition}

The main routes developed in this volume are cluster and phase-cell
expansions, which construct measures and Schwinger functions by estimates;
lattice constructions, which define finite systems exactly and study scaling
limits; and stochastic quantization or singular SPDE, which construct
invariant Euclidean measures through dynamics.
These routes overlap in low-dimensional scalar examples.  Their comparison is
part of the content: agreement of independent constructions requires a
separate identification theorem in the framework where the limiting objects
are compared.

\section{Status Catalog}

\begin{table}[H]
\centering
\scriptsize
\renewcommand{\arraystretch}{0.98}
\begin{tabular}{@{}>{\raggedright\arraybackslash}p{0.19\textwidth}
  >{\raggedright\arraybackslash}p{0.34\textwidth}
  >{\raggedright\arraybackslash}p{0.37\textwidth}@{}}
\hline
Theory or regime & Constructive status & Framework boundary \\
\hline
Free scalar, spinor, and generalized-free fields &
Constructed Wightman and Euclidean examples with explicit correlation
functions and positive Hilbert-space reconstruction. &
They calibrate the frameworks; interacting local dynamics require additional
construction. \\
\(P(\phi)_2\), including \(\phi^4_2\) &
Constructed non-Gaussian Euclidean and Wightman models for bounded-below
polynomial interactions after Wick ordering and local renormalization. &
Cluster expansion and Euclidean measure methods verify OS/Wightman
properties in two spacetime dimensions. \\
\(\phi^4_3\), massive regime &
Constructed non-Gaussian scalar models with mass, coupling, and vacuum-energy
renormalization. &
This is the fundamental three-dimensional scalar constructive example; the
continuum measure is singular with respect to the Gaussian reference. \\
Sine--Gordon in two dimensions &
Constructed in parameter ranges where the vertex interaction is
renormalizable and the infrared setup is controlled. &
The bosonization relation to massive Thirring-type descriptions requires its
own operator and Hilbert-space identifications. \\
Yukawa\(_2\) and related scalar--fermion models &
Constructive results exist for specific two-dimensional models after fixing
fermionic determinants or Pfaffians, Wick ordering, and ultraviolet
renormalization. &
Each theorem has model-dependent hypotheses; the Lagrangian label alone does
not define a completed construction. \\
Gross--Neveu\(_2\), massive asymptotically free regime &
Constructive renormalization-group constructions exist for massive
two-dimensional fermionic models. &
They are important tests of rigorous asymptotic freedom in a non-gauge
setting. \\
Two-dimensional Yang--Mills &
Constructed by heat-kernel, holonomy, and projective-limit methods for compact
gauge group. &
This is a gauge-theory construction with no local propagating gluons; it does
not settle four-dimensional Yang--Mills. \\
Stochastic quantization of \(\phi^4_2\) and \(\phi^4_3\) &
Constructed through Da Prato--Debussche, regularity structures,
paracontrolled distributions, and related singular-SPDE methods, with
renormalized invariant measures. &
The SPDE route gives a dynamic construction of Euclidean measures and a
precise algebraic meaning to the counterterms in the stochastic equation. \\
Three-dimensional Ising scaling limit as a CFT &
Open as a constructive conformal field theory with full operator algebra and
OS/Wightman reconstruction. &
Numerical and bootstrap data are extremely constraining, but a continuum CFT
construction remains a separate theorem. \\
\(\phi^4_D\), \(D\ge5\), standard positive lattice route &
Triviality theorems force Gaussian scaling limits in the standard critical
Ising/\(\lambda\phi^4\) family. &
This rules out that constructive route to an interacting scalar continuum
limit above four dimensions. \\
\(\phi^4_4\), standard critical Ising/\(\lambda\phi^4\) route &
Marginal triviality theorems give Gaussian scaling limits for the standard
reflection-positive critical problem. &
The logically different possibility of another four-dimensional UV-complete
QFT agreeing with formal \(\phi^4_4\) perturbation theory to all orders
remains a separate open question. \\
Four-dimensional pure Yang--Mills &
Finite Wilson lattices are rigorous compact-group integrals; continuum
existence with a mass gap is the Clay Millennium problem. &
The desired limit must construct gauge-invariant Schwinger functions or an
equivalent local observable theory satisfying positivity, locality, and
spectral properties. \\
QCD and the Standard Model &
Available nonperturbative definitions are hybrid: lattice gauge theory for
QCD-like sectors, perturbative and effective descriptions for electroweak and
heavy-scale observables, and experimentally fixed parameters. &
A single complete constructive continuum theorem for the full Standard Model
is presently an open construction. \\
\hline
\end{tabular}
\caption{Constructive status of representative QFT regimes.  The entries
separate completed constructions, triviality theorems, finite-regulator
definitions, and open continuum-limit problems.}
\label{tab:constructive-master-status}
\end{table}

\emph{Reference anchors for Table~\ref{tab:constructive-master-status}.}
The scalar constructive line runs through Nelson, Glimm--Jaffe,
Glimm--Jaffe--Spencer, Simon, Feldman--Osterwalder, Magnen--Seneor, Brydges,
Fr\"ohlich, and Sokal.  The two-dimensional fermionic and asymptotically free
examples include work of Feldman, Magnen, Rivasseau, and Seneor.  Rigorous
two-dimensional Yang--Mills may be approached through Driver, Sengupta, Levy,
and related heat-kernel/holonomy constructions.  The stochastic route uses
Da Prato--Debussche, Hairer, Gubinelli--Imkeller--Perkowski,
Kupiainen, Mourrat--Weber, and later singular-SPDE developments.  Scalar
triviality above and at four dimensions is represented by Aizenman,
Fr\"ohlich, and Aizenman--Duminil-Copin.

\section{Stochastic Quantization as a Constructive Route}

For Euclidean scalar \(\phi^4\), stochastic quantization starts from the
formal Langevin equation
\[
  \partial_t\Phi
  =
  \Delta\Phi-m^2\Phi-\lambda\Phi^3+\xi,
\]
where \(t\) is an auxiliary stochastic time and \(\xi\) is spacetime-time
white noise in the stochastic variables.  In \(D=2,3\), \(\Phi\) is a
distribution, so \(\Phi^3\) is defined only after renormalization.  The
renormalized equation has the form
\[
  \partial_t\Phi
  =
  \Delta\Phi-m^2\Phi-\lambda\Phi^3
  +c_1(\varepsilon)\Phi+c_0(\varepsilon)+\xi_\varepsilon
\]
before the cutoff \(\varepsilon\downarrow0\) is removed, with the precise
counterterms determined by the chosen model and symmetry.  Regularity
structures, paracontrolled distributions, and rigorous RG methods provide
the local expansion and reconstruction technology that make this limit a
theorem in the constructed cases.

\begin{openproblem}[Four-dimensional constructive gauge theory]
\label{op:constructive-four-dimensional-gauge}
Construct four-dimensional pure Yang--Mills, or QCD in a stated parameter
regime, as a continuum local QFT with gauge-invariant observables, reflection
positivity or Hilbert-space positivity, Euclidean and Lorentzian covariance,
locality, and the spectral properties needed for particle and flux-tube
questions.  Lattice finite-volume integrals, perturbative asymptotic freedom,
and numerical evidence are indispensable inputs, but the continuum existence
and mass-gap theorem remains open.
\end{openproblem}
